\section{Estrutura do documento}

% 1
O Capítulo~\ref{sec_Introducao_Dissertacao}: Introdução ao tema, incluindo contextualização, motivação, pergunta de pesquisa, objetivos e metodologia.

% 2
O Capítulo~\ref{sec_Revisao_bibliografica}: Revisão da literatura, abordando trabalhos relacionados à dinâmica orbital relativa, robótica espacial e modelagem de operações de RVD/B.

% 3
O Capítulo~\ref{sec_Procedimento_simulacao}: Procedimento de simulação, com fluxogramas e organização da implementação computacional.

% 4
O Capítulo~\ref{sec_Formulacao_Matematica}: Formulação matemática, incluindo modelagem orbital, cinemática do manipulador via DH, formulação lagrangiana e controle de atitude.

% 5
O Capítulo~\ref{sec_resultados}: Ambiente de simulação e resultados obtidos para as fases de aproximação de longa distância, curta distância ede proximidade (final)

% 6
O Capítulo~\ref{sec_Discussao}: Discussão dos resultados, destacando comparação com a literatura, contribuições do trabalho, lições aprendidas e limitações.

% 7
O Capítulo~\ref{sec_Conclusao}: Conclusão, incluindo o atingimento dos objetivos, resumo das contribuições, relevância e trabalhos futuros.